\chapter{Mixed Graphical Models}
Mixed graphical models (MGMs) refer to multivariate models which represent cross-sectional relationships among continuous, binary, categorical, and/or count variables in terms of an undirected graph. Vector autoregressive (VAR) models are multivariate models that can be used to estimate time-lagged relationships among sets of continuous variables. Similar to the MGM, the mVAR model has been proposed as an extension of VARs to accommodate mixed variable types, including binary, categorical, and count variables \cite{haslbeck2015mgm}. A recent extension of MGMs now affords the inclusion of moderating influences (i.e., interaction effects) in such models, making them powerful tools for detecting complex structure in high-dimensional datasets \cite{haslbeck2018moderated}. The goal of the present work is to propose this extension to mVAR models as well, specifying them in a way that makes it possible to estimate time-lagged relationships with moderator effects.

\cite{wainwright2007high} describe a detailed method for estimating high-dimensional graphical models using $\ell_1$-regularization, which I will describe in greater detail in this paper \cite{wainwright2007high}. 

A random variable $X$ is a measurable function $X:\Omega\rightarrow S$ from the sample space $\Omega$ to a measurable space $S$ called the \textit{state space}. If $A \subset S$, the notation Pr$(X\in A)$ is shorthand for P$(\{\omega\in\Omega : X(\omega)\in A\})$. Events $A$ and $B$ are conditionally independent given $C$ \textit{iff}, given the knowledge that $C$ occurs, knowledge of whether $A$ occurs gives no additional information on the likelihood that $B$ occurs, and knowledge of whether $B$ occurs provides no information on the likelihood of $A$ occurring.

\section{General Mixed Graphical Models}
Consider a $p$-dimensional random vector $X$ = $(X_1, \ldots, X_p)$ with each variable $X_s$ taking values in a potentially different set $\mathcal{X}_s$, and let $\G = (V,E)$ be an undirected graph over $p$ nodes corresponding to the $p$ variables. Now suppose the node-conditional distribution of node $X_s$ conditioned on all other variables $X_{\setminus s} := X_{V \setminus s}$ is given by an arbitrary univariate exponential family distribution
\begin{equation} \label{nodeCond}
    P(X_s | X_{\setminus s}) = \exp\{E_s (X_{\setminus s}) \phi_s (X_s) + B_s (X_s) - \Phi (X_{\setminus s})\}
\end{equation}
where the functions of the sufficient statistic $\phi_s (\cdot)$ and the base measure $B_s (\cdot)$ are specified by the choice of exponential family and the canonical parameter $E_s (X_{\setminus s})$ is a function of all variables except $X_s$ (see for details) \cite{wainwright2008graphical}.

These node-conditional distributions are consistent with a joint distribution over the random vector $X$ as in cliqueprod, that is Markov with respect to graph $\G = (V,E)$ with the set of cliques $\mathcal{C}_{k}$ of size at most $k$, iff the canonical parameters $\{E_s (\cdot)\}_{s \in V}$ are a linear combination of products of univariate sufficient statistic functions $\{\phi (X_r)\}_{r \in N(s)}$ of order up to $k$
\begin{equation} \label{suffStat}
    \theta_s + \sum_{r \in N(s)} \theta_{s,r} \phi_{r} (X_r) + \ldots + \sum_{r_1, \ldots , r_{k-1} \in N(s)} \theta_{r_1, \ldots , r_{k-1}} \prod_{j=1}^{k-1} \phi_{r_j} (X_{r_j})
\end{equation}
where $\theta_{s \cdot} \, := \, \{ \theta_s , \theta_{s,r} , \ldots , \theta_{s r_2 \ldots r_k}\}$ is a set of parameters and $N(s)$ is the set of neighbors of node $s$ according to graph $\G$ \cite{yang2014mixed}. Factorizing $p$ conditional distributions as in \eqref{nodeCond} gives the joint distribution
\begin{equation} \label{jointMGM}
    \begin{split}
        P(X) = \exp\Bigg\{&\sum_{s \in V} \theta_s \phi_s (X_s) + \sum_{s \in V} \sum_{r \in N(s)} \theta_{s,r} \phi_s (X_s) \phi_r (X_r) + \\
        &\cdots + \sum_{r_1 , \ldots , r_k \in \mathcal{C}} \theta_{r_1 , \ldots , r_k} \prod_{j=1}^{k} \phi_{r_j} (X_{r_j}) + \sum_{s \in V} B_s (X_s) - \Phi (\theta) \Bigg\}
    \end{split}
\end{equation}
The dimensionality of the parameter vector $\theta$ depends on both the type of modeled variable and the order of interactions. If only continuous variables with pairwise interactions are modeled $(k=2)$, the MGM simplifies to the multivariate Gaussian distribution which is parameterized by a $1 \times p$ vector of intercepts and a $p \times p$ matrix of $\binom{p}{2}$ partial correlations. Including all 3-way interactions would lead to an additional $\binom{p}{3}$ parameters, etc. 

Necessary conditions for the mixed density in \eqref{jointMGM} are discussed in \cite{yang2014mixed}. \cite{chen2014selection} show constraints on the parameter space to ensure normalizability for a number of MGMs with at most pairwise interactions. 

\subsection{The Ising-Gaussian Model}
The Ising-Gaussian model would be a specific example of the joint distribution in equation \ref{jointMGM}. Consider a random vector $X := (Y,Z)$, where $Y = \{Y_1,\ldots,Y_{p_1}\}$ are univariate Gaussian random variables, $Z = \{Z_1,\ldots,Z_{p_2}\}$ are univariate Bernoulli random variables and we only consider pairwise interactions between sufficient statistics. For the univariate Gaussian distribution (with known variance $\sigma^2$) the sufficient statistic function is $\phi_Y (Y_s) = \frac{Y_s}{\sigma_s}$ and the base measure is $B_Y (Y_s) = -\frac{Y_{s}^2}{2\sigma_{s}^2}$. The Bernoulli distribution has the sufficient statistic function $\phi_{Z_s} = Z_s$ and the base measure $B_Z(Z_s)=0$. From the MGM joint distribution in equation \ref{jointMGM} follows that this mixed density is given by
\begin{equation} \label{jointIG}
    \begin{split}
        P(Y,Z) \propto \exp\Bigg\{&\sum_{s \in V_Y} \frac{\theta_s}{\sigma_s} Y_s + \sum_{r \in V_Z} \theta_r Z_r + \sum_{(s,r) \in E_Y} \frac{\theta_{s,r}}{\sigma_s\sigma_r} Y_s Y_r + \\
        &\sum_{(s,r) \in E_Z} \theta_{s,r} Z_s Z_r + \sum_{(s,r) \in E_{Y Z}} \frac{\theta_{s,r}}{\sigma_s} Y_s Z_r - \sum_{s \in V_Y} \frac{Y_{s}^2}{2\sigma_{s}^2} \Bigg\}
    \end{split}
\end{equation}
where the first two terms are thresholds for Gaussian and Bernoulli variables, the third term represents pairwise interactions between Gaussians, the fourth term represents pairwise interactions between Bernoulli variables, the fifth term represents pairwise interactions between Gaussians and Bernoulli variables, and the last term sums over the base measures for the Gaussians.

When the conditional distribution is a Bernoulli random variable $Z_r$, it is given by
\begin{equation} \label{bernCond}
    P(Z_r|Z_{\setminus r},Y) \propto \exp \Bigg\{\theta_r Z_r + \sum_{s \in N(r)_{Z}} \theta_{s,r} Z_s Z_r + \sum_{s \in N(r)_{Y}} \frac{\theta_{s,r}}{\sigma_r} Z_r Y_s \Bigg\}
\end{equation}
The conditional distribution in equation \ref{bernCond} has the same form as the distribution of a single variable conditioned on all remaining variables in an Ising model plus one additional term for interactions between Bernoulli and Gaussian random variables.

When the conditional distribution is a Gaussian random variable $Y_s$, it is given by
\begin{equation} \label{gaussCond}
    P(Y_s|Y_{\setminus s},Z) \propto \exp \Bigg\{\frac{\theta_s}{\sigma_s} Y_s + \sum_{r \in N(s)_{Y}} \frac{\theta_{s,r}}{\sigma_s \sigma_r} Y_s Y_r + \sum_{r \in N(s)_{Z}} \frac{\theta_{s,r}}{\sigma_s} Y_s Z_r - \frac{Y_{s}^{2}}{2\sigma_{s}^{2}} \Bigg\}
\end{equation}
Let $\sigma=1$, factor out $Y_s$ and let $\mu_s = \theta_s + \sum_{r \in N(s)_{Y}} \theta_{s,r} Y_r + \sum_{r \in N(s)_{Z}} \theta_{s,r} Z_r$. Finally, when taking $\frac{\mu_{s}^{2}}{2}$ out of the log normalization constant, we arrive with basic algebra at the well-known form of the univariate Gaussian distribution with unit variance
\begin{equation} \label{uniGauss}
    P(Y_s|Y_{\setminus s},Z) = \frac{1}{\sqrt{2 \pi}} \exp \bigg\{-\frac{(Y_s - \mu_s)^2}{2} \bigg\}
\end{equation}


\subparagraph{Relationship between model parameters and edges in graph}
For pairwise MGMs (size of cliques is at most $k=2$), a pairwise interaction between two continuous variables $s$ and $r$ is parameterized by a single parameter $\theta_{s,r}$. Thus, whether $s \sim r$ holds depends upon whether $\theta_{s,r}$ is zero or not; i.e., $(s,r) \in E \implies \theta_{s,r} \neq 0$. Therefore if only pairwise interactions between continuous variables are modeled, any given edge is a function of a single parameter. Interactions between categorical variables with $m > 2$, however, are specified by more than one parameter. For example, an interaction between two categorical variables with $m$ and $u$ categories is parameterized by $R=(m-1) \times (u - 1)$ parameters associated with the corresponding indicator functions for all $R$ states \cite{agresti2003cat}. A pairwise interaction between a categorical variable with $m$ categories and a continuous variable has $R=1 \times (m-1)$ parameters associated with $m-1$ indicator functions multipled with the continuous variable. In this case $\theta_{s,r}^{z}$ is a parameter defining the interactions between the nodes $s$ and $r$ indexed by $z \in \{1,\ldots,R\}$. 

In the \code{mgm} package the parameterization for multinomial regression of \code{glmnet} is used to model the probability of each state of the predicted variable, coding the first category of the predictor variable as the reference category that is absorbed in the intercept \cite{friedman2010bregularization}. This results in $m \times (u-1)$ parameters, where $m$ is the number of categories in the predicted variable. Here, an edge is present if \textit{any} parameters in $\theta_{s,r}$ is nonzero.

\section{Estimating Mixed Graphical Models}
The joint distribution in equation \ref{jointMGM} can be represented as a factorization of univariate conditional distributions. Thus, if we estimate $p$ univariate conditional distributions with the parameterization in equation \ref{suffStat}, we obtain the joint distribution. Given that all univariate conditional distributions are members of the univariate conditional distributions are members of the exponential family, it is possible to estimate the joint distribution in equation \ref{jointMGM} by a series of $p$ regressions in the Generalized Linear Model (GLM) framework \cite{nelder1972generalized}. From a graphical modeling perspective, this means that we estimate the neighborhood $N(s)$ of each node $s \in V$ and then combine all neighborhoods to obtain an estimate of the graph $\G$ \cite{meinshausen2006high}.

In order to obtain parameter estimates that are exactly zero, we minimize the negative log-likelihood $\L(\theta,X)$ together with the $\ell_1$-norm of the parameter vector
\begin{equation} \label{lassoMGM}
    \hat{\theta}=\argmin_{\theta} \{\L(\theta,X) + \lambda_n \norm{\theta}_{1}\}
\end{equation}
where $\norm{\theta}_{1} = \sum_{j=1}^{J} |\theta_j|$ and $J$ is the length of the parameter vector $\theta$. The negative log-likelihood $\L(\theta,X)$ is defined as the exponential family distribution of the node under consideration. For Gaussian variables, minimizing the log-likelihood is equivalent to minimizing the squared loss $\L(\theta,X) = \lVert X_s - X_{\setminus s} \theta \rVert_{2}^{2}$. Simply put, this is performing an $\ell_1$-penalized (LASSO) regression in the GLM framework with a link-function appropriate to the node under consideration \cite{nelder1972generalized}. The reason for using the $\ell_1$-penalty is to ensure that the model is identified in high-dimensional settings where $p > n$, where we have more parameters than observations \cite{hastie2015statistical}. 

The design matrix is defined with respect to the conditional distribution of node $s$ in the $k$-order MGM. Thus, if $k=2$, the design matrix for the regression on node $s$ contains all other variables or the corresponding indicator functions (for categorical variables). If $k=3$, the design matrix for the regression on node $s$ contains all other variables or the corresponding indicator functions, plus the products of all pairs of variables in $V_{\setminus s}$, or the $(m-1) \times (u-1)$ indicator functions in the case of categorical variables with $m$ and $u$ categories.

To give non-asymptotic guarantees of false and true positive rates for the $\ell_1$-regularized regression estimator it is necessary to put a lower bound $\tau_n$ on the size of the parameters in the true model. This assumption is referred to as the \textit{beta min condition} (see e.g.) \cite{hastie2015statistical}. By thresholding the estimates at $\tau_n$, we approximately enforce this condition (see also) \cite{loh2013structure}. For estimating the joint distribution in equation \ref{jointMGM}, it is shown in \cite{haslbeck2017structure} that 
\begin{equation} \label{threshHW}
    \tau_n \asymp s_0 \sqrt{\frac{\log p}{n}} \leq s_0 \lambda_n
\end{equation}
where $s_0$ is the true number of neighbors. If all variables are continuous, the number of neighbors is equal to the number of nonzero parameter estimates $s_0 = \norm{\theta^{*}}_0$, where $\theta^{*}$ is the true parameter vector. With categorical variables, given that there are several parameters, the categorical neighbor is present if at least one of the parameters defining the interaction is nonzero. Since $\theta^{*}$ is unknown, we use the estimated parameter $\hat{\theta}$ to obtain the \textit{estimated} number of neighbors $\hat{s}_0 = \norm{\hat{\theta}}_0$. For interactions involving more than one parameter, we plug in the aggregated parameter (AND/OR rule). Remember again that the thresholding parameter $\tau_n$ guarantees false and true positive rates.

The above is what was used by \cite{haslbeck2015mgm}. However, \cite{haslbeck2017structure} use a different formulation (as well as in their \code{mgm} function), which is informed by the results of \cite{loh2013structure}:
\begin{equation} \label{threshLW}
    \tau_n \asymp \sqrt{d} \norm{\theta^*}_2 \sqrt{\frac{\log p}{n}}
\end{equation}
Note that we use $\hat{\theta}$ to stand in for $\theta^*$. Also, $p$ refers to the number of predictors for the particular nodewise regression at hand (which may differ from $p$ of the dataset, if categorical variables are present and broken up into dummy-coded variables), and $d$ refers to the true degree of the graph, which is equal to the order minus one $(k-1)$.

In addition to determining whether or not each edge is nonzero, the \textit{weight} is computed from the set of parameters of each interaction. If the interaction only involves continuous variables, there is only one parameter to take into account. If there are categorical variables, however, we take the mean of the absolute value of all parameters as the weight of the edge. From the nodewise regressions we obtain $k$ edge-weights for each $k$-order interaction. For instance, for the pairwise interaction $(k=2)$ between nodes $s$ and $r$, we obtain one parameter $\theta_{s,r}$ from the regression on $s$, and $\theta_{r,s}$ from the regression on $r$. These will then be combined based on either the OR-rule (simply take the arithmetic mean of $k$ parameter estimates) or the AND-rule (take the arithmetic mean of $k$ parameter estimates if all parameter estimates are nonzero, otherwise set the parameter to zero). 

\subparagraph{Estimating mixed graphical models via neighborhood regression}
For each $s \in V$
\begin{enumerate}
    \item Construct design matrix defined by $k$, the order of the MGM
    \item Solve the lasso problem \eqref{lassoMGM} with regularization parameter $\lambda_n$; this can be done via goodness-of-fit indices or cross-validation
    \item Threshold the estimates at $\tau_n$, using either \eqref{threshLW} or \eqref{threshHW}
    \item Aggregate interactions with several parameters (if there are categorical variables) into a single edge-weight
\end{enumerate}
After this is complete, (a) combine the edge-weights across the nodewise regressions using either the AND- or OR-rule, and (b) define $\G$ based on the zero/nonzero pattern in the combined parameter vector.

As mentioned above, the regularization parameter $\lambda_n$ can be selected using either cross-validation or a model-selection criterion such as the Extended Bayesian Information Criterion (EBIC):
\begin{equation} \label{ebic}
    EBIC_{\gamma} (\hat{\theta}) = -2\L(\hat{\theta}) + \hat{s}_0 \log n + 2 \gamma \hat{s}_0 \log p
\end{equation}
where $\L$ is the log likelihood of the conditional density specified by the estimated parameter vector $\hat{\theta}$, $\hat{s}_0$ is the number of nonzero neighbors in the candidate model, and $\gamma$ is the tuning parameter. If $\gamma = 0$ then the EBIC is identical to the BIC \cite{schwarz1978estimating}. The EBIC has been shown to perform well asymptotically in selecting sparse graphs \cite{foygel2010extended,foygel2015high}. \cite{foygel2010extended} used values $\gamma \in \{0,.25,.5,.75,1\}$ and showed that increasing $\gamma$ from 0 to 0.25 led to considerable decrease in false positives without leading to much of an increase in false negatives. Thus, setting $\gamma = 0.25$ might be an appropriate rule of thumb, although a simulation study could be conducted based on your own data in order to determine the \textit{optimal} value of $\gamma$. This can be done using the \code{mgm} package in \code{R} \cite{haslbeck2015mgm}. The computational complexity of the neighborhood regression procedure is quite poor, $\mathcal{O}(p\log(p2^{k-1}))$, and so the method does not scale well for large $k$.

\section{Mixed Vector Autoregressive Models (mVAR)}
For VAR models, each node $s$ at time $t$ is modeled as a linear combination of all variables (including $s$) given some lag $j$. The standard VAR is defined as a Gaussian white noise process, such that the model can be split into $p$ conditional Gaussian distributions \cite{hamilton1994time,pfaff2008time}. Instead of a univariate conditional Gaussian distribution, one can associate other univariate exponential family members with a given node. This produces almost the exact same model and scenario as is present for the MGM estimation method. The only difference is that the canonical parameter of the node-conditional is not a function of parameters associated with interactions at the \textit{same time point}, but rather at some previous \textit{lagged time point(s)}. 

The mVAR can be estimated by estimating the parameters of the conditional probability of each variable $s$ as a function of all variables (including itself) at a set of specified previous time points, denoted by $L$. For instance, $L=\{1,2,3\}$ specifies a VAR model with lags 1, 2, and 3. The time index is here treated as a superscript $t$ for all variables since we now have time ordered observations. We must now define the canonical parameter $E_{s}^{t}(X)$ of the conditional distribution $P(X_{s}^{t}|X^{t-1},X^{t-2},X^{t-3})$ at time $t$
\begin{equation} \label{mVAR}
    E_{s}^{t}(X) = \theta_s + \sum_{j \in L} \sum_{r \in N(s)} \theta_{s,r}^{t-j} \phi_r (X_{r}^{t-j})
\end{equation}
Only pairwise interactions are included here because \code{mgm} does not implement higher order interactions for mVAR models (yet!). The canonical parameter function in equation \ref{mVAR} defines the negative log-likelihood $\L(\theta,X)$ in equation \ref{lassoMGM} and thus we can use the primary algorithm with two modifications: (a) we define the design matrix as a function of lags $L$ instead of maximal order of interactions (which are here fixed to $k=2$, and (b) we do not apply the AND/OR rule, because the cross-lagged effect of $X_{s}^{t-1}$ on $X_{r}^{t}$ is a different effect than the cross-lagged effect of $X_{r}^{t-1}$ on $X_{s}^{t}$ and so no parameter is estimated twice. Here is the algorithm
\begin{enumerate}
    \item For each $s \in V$
    \begin{enumerate}
        \item Construct design matrix defined by $L$
        \item Solve the lasso problem \eqref{lassoMGM} with regularization parameter $\lambda_n$
    \end{enumerate}
    \item Threshold the estimates at $\tau_n$
    \item Define the directed graph $D_j$ based on the zero/nonzero pattern in the combined parameter vector for each lag $j \in L$
\end{enumerate}
The computational complexity here is $\mathcal{O}(p\log(p|L|))$. The directed graphs in the $p \times p \times |L|$ array $D$ are not encoding conditional independence statements as in the graph $\G$ for MGMs. \cite{haslbeck2017estimate} show that this method can recover VAR models with Gaussian noise in a variety of situations.

\subsection{Time-varying mVAR models}
For mVAR models, it is assumed that the models are stationary, where observations at each time $t$ are generated from the same distribution parameterized by $\theta$. With time-varying models, the assumption is relaxed such that parameters $\theta^t$ can be different at each time point $t \in \mathcal{T}_n = \{\frac{1}{n},\frac{2}{n},\ldots,1\}$, where $n$ is the number of time points in the series. 

We make assumptions about how the parameters of the true model vary as a function of time. These are usually assumptions about \textit{local stationarity} (e.g.,) \cite{zhou2010time}, which comes in two flavors: (a) we assume there exists a partition $\mathcal{B}$ of $\mathcal{T}_n$ in which time points are consecutive and in each of the subsets $B \in \mathcal{B}$ the model is stationary, i.e., $\forall i,j \in B : \theta^i = \theta^j$. These piecewise constant time-varying models can be estimated with a fused lasso penalty, which puts an additional penalty on parameter changes from one time to the next, or (b) we require that the model $\theta^t$ is a smooth function of time. In this case we combine observations close in time for estimation, because we know their generating models are similar. To do this we fit local models $\hat{\theta}^{t^e}$ across the time series, which only give high weight to data points close to the given estimation point $t^e$. The full time-varying model is then the set of all the local estimates $\{\theta^{e_1},\theta^{e_2},\ldots,\theta^{|\mathcal{E}|}\}$ at estimation points $\mathcal{E} = \{t_1^e,t_2^e,\ldots,t^{|\mathcal{E}|}\}$, where the entries in $\mathcal{E}$ are usually equally spaced across the time series and the number of estimation points $|\mathcal{E}|$ is chosen depending on how fine-grained one would like to describe $\theta^t$ as a function of time $t$.

Stating this formally, we estimate the model $\theta^{t^e}$ at time point $t^e$ by minimizing a weighted version of the loss function in equation \ref{lassoMGM}
\begin{equation} \label{lassoTV}
    \hat{\theta}^{t^e} = \argmin_{\theta} \Bigg\{\frac{1}{\sum_{t=1}^{n} w_t^{t^e}} \sum_{t=1}^n w_t^{t^e} \L(\theta,X^t) + \lambda_n \norm{\theta}_1 \Bigg\}
\end{equation}
where $w_t^{t^e}$ as a function of $t$ is defined by a kernel centered over $t^e$. Specifically, we define the weight function $w_t^{t^e}$ to be a Gaussian kernal, normalized so that the largest weight is equal to 1 \cite{zhou2010time}
\begin{equation} \label{weight}
    w_t^{t^e} = \frac{Z_t}{\max_{(t\in\mathcal{T}_n)} \{\bigcup_t Z_t\}}, \quad \text{where} \quad Z_t = \frac{1}{\sqrt{2\pi}} \exp \Bigg\{-\frac{(t-t^e)^2}{2\sigma^2}\Bigg\}
\end{equation}
The sum of all weights $n_{\sigma,t^e} = \sum_{t=1}^n w_t^{t^e}$ (or area under the curve) used at a given estimation point $t^e$ indicates the amount of data used for estimation at $t^e$ relative to the full series.